\documentclass[UTF8,a4paper,11pt]{ctexart}

\usepackage{amsmath,amssymb}
\usepackage{geometry}
\geometry{margin=2.5cm}
\usepackage{booktabs}
\usepackage{hyperref}
\usepackage{enumitem}

\title{时域与频域：周期 $T$ 的含义}
\author{}
\date{}

\begin{document}
\maketitle

在这张图里，\textbf{$T_1$ 和 $T_2$ 代表两个不同正弦波的周期（Period）}。

左下角画的是：

\[
A_1\sin\left(\frac{2\pi}{T_1}t\right)
\]

和

\[
A_2\sin\left(\frac{2\pi}{T_2}t\right)
\]

它们相加后形成中间那个更复杂的波形。

\bigskip\noindent\rule{\linewidth}{0.4pt}\bigskip

\subsection{什么是周期 $T$？}

周期表示：

\begin{quote}
\textbf{完成一次完整振动所需要的时间。}
\end{quote}

例如：

\begin{itemize}
  \item $T=1\,\mathrm{s}$ $\rightarrow$ 每 1 秒振动一次
  \item $T=0.1\,\mathrm{s}$ $\rightarrow$ 每 0.1 秒振动一次
\end{itemize}

周期越小：

\begin{itemize}
  \item 振动越快
  \item 频率越高
\end{itemize}

\bigskip\noindent\rule{\linewidth}{0.4pt}\bigskip

\subsection{频率和周期的关系}

\[
f=\frac{1}{T}
\]

例如：

\begin{center}
\begin{tabular}{cc}
  \toprule
  周期 $T$ & 频率 $f$ \\
  \midrule
  1 s & 1 Hz \\
  0.5 s & 2 Hz \\
  0.1 s & 10 Hz \\
  \bottomrule
\end{tabular}
\end{center}

\bigskip\noindent\rule{\linewidth}{0.4pt}\bigskip

\subsection{图中为什么写 $1/T_1$ 和 $1/T_2$？}

右边其实是在引入\textbf{频域（Frequency Domain）}。

时域中：

\begin{itemize}
  \item 黄色大波：振幅 $A_1$，周期 $T_1$
  \item 黄色小波：振幅 $A_2$，周期 $T_2$
\end{itemize}

相加后得到复杂波形。

频域中：

\begin{itemize}
  \item 在频率 $1/T_1$ 处有一个峰值 $A_1$
  \item 在频率 $1/T_2$ 处有一个峰值 $A_2$
\end{itemize}

意思是：

\begin{quote}
这个复杂波形实际上由两个频率成分组成。
\end{quote}

\bigskip\noindent\rule{\linewidth}{0.4pt}\bigskip

这正是控制工程里从\textbf{时域（Time Domain）}到\textbf{频域（Frequency Domain）}的核心思想。

时域看到的是：

\begin{itemize}
  \item 一条复杂曲线
\end{itemize}

频域看到的是：

\begin{itemize}
  \item 「里面包含哪些振动频率，以及每个频率有多强」
\end{itemize}

就像音乐：

\begin{itemize}
  \item 时域：麦克风录到的声音波形
  \item 频域：低音 100 Hz、中音 1000 Hz、高音 5000 Hz 各有多少能量
\end{itemize}

图里的 $T_1$、$T_2$ 就是在为后面要学的 \textbf{傅里叶变换（Fourier Transform）} 和 \textbf{Bode Plot（波特图）} 做铺垫。控制系统中的 $G(s)$ 最终经常就是在频域里分析的。

\end{document}
